\subsection{完全切割族}

设 A 是一个环，M 是一个 A-模， \(x = (x_{i})_{i \in I}\) 是 A 中元素的一个族。

定义 2. --- 称族 x 对 M 是完全切割的，如果有 \(\mathrm{H}_{i}(\mathbf{x},\mathbf{M})=0\) 对 i\textgreater0 。

若 I 是有限的，则这等价于（X，第 150 页）说有 \(\mathbf{H}^i (\pmb {x},\mathbf{M}) = 0\) 对 \(i <   \mathrm{Card}\) (I)。

下面的命题使我们能够给出完全切割族的一些例子。

命题 5. --- 设 \(x = (x_1, \ldots, x_n)\) 是 A 中元素的一个序列。如果对 \(i=1, \ldots, n\) ，比率为 \(x_i\) 的位似在 A-模 \(M/(x_1 M + \cdots + x_{i-1} M)\) 中是单射的，则序列 \(x\) 对 M 是完全切割的。

满足命题条件的序列称为对 M 正则的，或 M-正则的。注意，此性质一般依赖于 \(x_i\) 的顺序。例如，序列 (1, 0) 总是 M-正则的，而序列 (0, 1) 仅当 M 为零模时才是 M-正则的。另一方面，一个序列是否完全切割并不依赖于各项的顺序。

让我们对 \(n\) 作归纳来证明该命题， \(n=0\) 的情形是显然的。设 \(x'=(x_{1},...,x_{n-1})\) ；若序列 \(x\) 是 M-正则的，则序列 \(x'\) 是 M-正则的，且以 \(x_{n}\) 乘 \(\mathrm{M}/(\mathbf{x}')\mathrm{M}\) 的位似是单射；由归纳假设，我们有 \(\mathrm{H}_{i}(\mathbf{x}',\mathrm{M})=0\) 对 \(i>0\) ；于是由 X，第 157 页，推论 3，可知 \(\mathrm{H}_{i}(\mathbf{x},\mathrm{M})=0\) 对 \(i>0\) 。

例. --- 设 \(k\) 是一个环；取 \(\mathbf{A} = k[\mathbf{X}_1, \ldots, \mathbf{X}_n]\) 和 \(\boldsymbol{x} = (\mathbf{X}_1, \ldots, \mathbf{X}_n)\) 。序列 \(\boldsymbol{x}\) 是 \(\mathbf{A}\) -正则的，且该命题重新给出次数 \(>0\) 时复形 \(\mathbf{S}_k(k^n) \otimes_k \boldsymbol{\Lambda}_k(k^n)\) 的零调性（参见 X，第 152 页，命题 3）。

同样地，在形式幂级数环 \(\hat{\mathbf{A}} = k[[\mathbf{X}_1, \ldots, \mathbf{X}_n]]\) 中，序列 \((\mathbf{X}_1, \ldots, \mathbf{X}_n)\) 是 \(\hat{\mathbf{A}}\) -正则的，因此对 \(\hat{\mathbf{A}}\) 完全切割。

命题 6. --- a) 若 \(\sum_{i\in I}x_i\mathbf{A} = \mathbf{A}\) ，则族 \((x_i)_{i\in I}\) 对 M 是完全切割的。

\begin{enumerate}
\def\labelenumi{\alph{enumi})}
\setcounter{enumi}{1}
\tightlist
\item
  设 \(x = (x_1, \ldots, x_n)\) 是 A 的元素序列。设 \((a_{ij}) \in \mathbf{GL}_n(\mathbf{A})\) ；令
\end{enumerate}

\[
y _ {i} = \sum_ {j} a _ {i j} x _ {j}.
\]

若序列 x 对 M 完全切割，则序列 \((y_{1}, \ldots, y_{n})\) 亦然。c) 设 \(k_{1}, \ldots, k_{n}\) 为整数 \(\geqslant 1\) ；序列 \((x_{1}^{k_{1}}, \ldots, x_{n}^{k_{n}})\) 对 M 完全切割，当且仅当序列 x 对 M 完全切割。断言 a) 由第 148 页推论 1 得出。

我们来证明 \(b\) )。设 \(f: \mathbf{A}^n \to \mathbf{A}^n\) 为由矩阵 \(^t(a_{ij})\) 定义的同构；由 X，第 151 页可知 \(1_{\mathbf{M}} \otimes \Lambda(f)\) 是复形 \(\mathbf{K}.(y, \mathbf{M})\) 到复形 \(\mathbf{K}.(x, \mathbf{M})\) 上的同构，从而得 \(b\) 。

为了证明 \(c\) ），显然只需证明：当 \(k\) 是整数 \(\geqslant 1\) 时，序列 \((x_{1}, \ldots, x_{n-1}, x_{n}^{k})\) 对 M 完全切割当且仅当序列 \(x\) 对 M 完全切割。设 \(x' = (x_{1}, \ldots, x_{n-1})\) ；由第 157 页推论 3，第一个条件（相应地，第二个条件）意味着比率为 \(x_{n}^{k}\) （相应地， \(x_{n}\) ）的位似在 H \(_{i}(x', M)\) 中对 \(i \geqslant 1\) 是双射的，且在 M/(x') M 中是单射的。这两个条件显然等价，因此得 \(c\) ）。

注记。--- 1) 对正则序列，与 c) 类似的断言成立（X，第 207 页，习题 5）。

命题 7. --- a) 设 N 是一个平坦的 A-模。如果族 x 对 M 完全切割，则它对 M ⊗\_A N 也完全切割。

\begin{enumerate}
\def\labelenumi{\alph{enumi})}
\setcounter{enumi}{1}
\tightlist
\item
  设 \(0 \to \mathbf{M}' \to \mathbf{M} \to \mathbf{M}'' \to 0\) 是 A-模的一个正合列。如果族 \(x\) 对 \(\mathbf{M}'\) 和 \(\mathbf{M}''\) 都完全切割，则它对 \(\mathbf{M}\) 也完全切割。
\end{enumerate}

复形 \(\mathbf{K}.\left(\mathbf{x},\mathbf{M}\otimes_{\mathbf{A}}\mathbf{N}\right)\) 由定义同构于 \(\mathbf{K}.\left(\mathbf{x},\mathbf{M}\right)\otimes_{\mathbf{A}}\mathbf{N}\) ；由于 N 是平坦的，我们推出一个同构 H. \((\mathbf{x},\mathbf{M})\otimes_{\mathbf{A}}\mathbf{N}\to\mathrm{H}.\left(\mathbf{x},\mathbf{M}\otimes_{\mathbf{A}}\mathbf{N}\right)\) （X，第 66 页，推论 2），由此得 \(a\) ）。

由于复形 \(\mathbf{K}.\left(\mathbf{x}\right)\) 是平坦的，我们有一个复形的正合列

\[
0 \rightarrow \mathbf {K}. (x, \mathrm{M} ^ {\prime}) \rightarrow \mathbf {K}. (x, \mathrm{M}) \rightarrow \mathbf {K}. (x, \mathrm{M} ^ {\prime \prime}) \rightarrow 0;
\]

断言 b) 由相伴的同调正合列得出。

注记。--- 2) 对正则序列而言，与 a) 和 b) 类似的断言是直接的。

\begin{enumerate}
\def\labelenumi{\arabic{enumi})}
\setcounter{enumi}{2}
\tightlist
\item
  如果族 \(x\) 对 \(A\) 完全切割，则复形 \(K.(x, A)\) 定义了 \(A\) -模 \(A/x\) 的一个自由分解，其中 \(x = \sum_{i \in I} x_i A\) ；因此对每个整数 \(i \geqslant 0\) 以及每个 \(A\) -模 \(M\) 我们有同构
\end{enumerate}

\[
\operatorname{Ext} _ {\mathbf {A}} ^ {n - i} (\mathbf {A} / \mathfrak {x}, \mathbf {M}) \rightarrow \mathrm{H} ^ {n - i} (\mathfrak {x}, \mathbf {M}) \rightarrow \mathrm{H} _ {i} (\mathfrak {x}, \mathbf {M}) \rightarrow \operatorname{Tor} _ {i} ^ {\mathbf {A}} (\mathbf {A} / \mathfrak {x}, \mathbf {M}).\tag{15}
\]

\begin{enumerate}
\def\labelenumi{\arabic{enumi})}
\setcounter{enumi}{3}
\tightlist
\item
  我们称序列 \(\boldsymbol{x} = (x_1, \dots, x_n)\) 是 M-余正则的，如果（记 \((x_i)_M\) 为 M 中比率为 \(x_i\) 的位似）模中比率为 \(x_i\) 的位似
\end{enumerate}

\[
\operatorname{Ker} \left(x _ {1}\right) _ {\mathbf {M}} \cap \dots \cap \operatorname{Ker} \left(x _ {i - 1}\right) _ {\mathbf {M}}
\]

对 \(i = 1, \ldots, n\) 是满射的。于是我们有 \(\mathbf{H}_{i}(\mathbf{x}, \mathbf{M}) = 0\) 对 \(i < n\) ；其证明方式与命题 5 相同。

\[
\mathrm{A} = k \left[ \mathrm{D} _ {1}, \dots , \mathrm{D} _ {n} \right],
\]

\[
\mathbf {Q} \text {-algebra},
\]

\(M = k[T_1, \ldots, T_n]\) ，赋予其一个 A-模结构，使得 \(D_i P = \partial P / \partial T_i\) 对所有 \(P \in M (1 \leqslant i \leqslant n)\) 成立。立即可验证序列 \((D_1, \ldots, D_n)\) 是 M-余正则的；结合第 155 页的例子，可以推出 \(k[T_1, \ldots, T_n]\) 在 \(k\) 上的德拉姆复形在次数 \textgreater{} 0 时是零调的。

